% PARTES ANTES DO TEXTO
% Uma linha iniciada por % está desativada. Remova apenas esse % para ativá-la.
%
% A ordem abaixo é a do Manual de Normalização do IFPI 2024 (seção 5.2.1).
% A seleção ativa reproduz o modelo de ARTIGO (seção 7.2), em que são
% obrigatórios apenas capa, folha de rosto, título, autor, resumo no idioma
% do documento, resumo em outro idioma e data de aprovação.
%
% Para monografia, dissertação ou tese (seção 5), passam a ser obrigatórios
% também a folha de aprovação, a ficha catalográfica e o sumário.
%
% Os itens opcionais estão prontos para edição em pre-textual/opcionais/.
% Os elementos pré-textuais são CONTADOS, mas não numerados (seção 3.5).

\imprimircapa
\imprimirfolhaderosto*

% A ficha catalográfica vai no verso da folha de rosto e é feita por
% bibliotecário. No artigo não é exigida; veja o arquivo antes de ativar.
% % A ficha catalográfica é feita por bibliotecário, com o número de registro
% no CRB, e vai no VERSO da folha de rosto (manual IFPI 2024, seção 5.2.1.2).
% Ela é obrigatória para monografia, dissertação e tese. No ARTIGO como TCC
% -- que é a modalidade deste modelo -- NÃO é exigida (seção 7.2).
%
% Primeiro, adicione a ficha oficial como pre-textual/ficha-catalografica.pdf.
% Depois, ative a inclusão deste arquivo em estrutura/pre-textual.tex.
% Não substitua a ficha oficial por dados bibliográficos inventados.
\includepdf[pages=-]{pre-textual/ficha-catalografica.pdf}


% A errata, se houver, vem logo após a folha de rosto.
% % Ative em estrutura/pre-textual.tex somente se houver uma errata.
\clearpage
\begin{errata}
\imprimirautor. \textbf{\imprimirtitulo}. \imprimirlocal, \imprimirdata.

\begin{center}
\begin{tabular}{llll}
  \toprule
  \textbf{Folha} & \textbf{Linha} & \textbf{Onde se lê} & \textbf{Leia-se} \\
  \midrule
  1 & 10 & termo incorreto & termo correto \\
  \bottomrule
\end{tabular}
\end{center}
\end{errata}
\clearpage


% Escolha somente UMA folha de aprovação, se exigida pelo curso.
% Obrigatória em monografia, dissertação e tese (seção 5.1.2.4).
% \imprimirfolhadeaprovacao
% \imprimirfolhadeaprovacaoduascolunas

% % Ative este arquivo em estrutura/pre-textual.tex e substitua o exemplo.
% Elemento opcional. Manual IFPI 2024, seção 5.1.2: a dedicatória fica
% alinhada à direita e ao final da folha.
\clearpage
\begin{dedicatoria}
	% Posiciona a dedicatória na parte inferior da página.
	\vspace*{\fill}
    
    % Alinhamento à direita.
    \begin{flushright}
    	Dedico este trabalho a todos que acreditaram que ele sairia.
    \end{flushright}
    
\end{dedicatoria}
\clearpage

% % Ative este arquivo em estrutura/pre-textual.tex e substitua o exemplo.
% Elemento opcional (manual IFPI 2024, seção 5.1.2). Bolsistas da Capes
% devem incluir o texto exigido pela Portaria nº 206/2018: "O presente
% trabalho foi realizado com apoio da Coordenação de Aperfeiçoamento de
% Pessoal de Nível Superior - Brasil (CAPES)".
\clearpage
\begin{agradecimentos}
Agradeço às pessoas e instituições que contribuíram para a realização deste
trabalho. Descreva aqui as contribuições que deseja reconhecer.
\end{agradecimentos}
\clearpage

% % Ative em estrutura/pre-textual.tex. Use uma frase com autoria conferida.
% Elemento opcional. Manual IFPI 2024, seção 5.1.2: citação curta, colocada
% depois dos agradecimentos; pode ser usada também na abertura das seções
% primárias.
\clearpage
\begin{epigrafe}
\vspace*{\fill}
\begin{flushright}
\emph{Insira aqui a frase escolhida.}\\
Nome de quem escreveu a frase
\end{flushright}
\end{epigrafe}
\clearpage


% Resumo no idioma do documento e em língua estrangeira: os dois são
% obrigatórios (seções 4 e 7.2).
% RESUMO: substitua o parágrafo de orientação pelo resumo da sua pesquisa.
% Título, nomes, notas e palavras-chave vêm de configuracoes/metadados.tex.
%
% Regras do manual IFPI 2024 (seções 4 e 7.4) para o resumo do artigo:
%   - parágrafo ÚNICO, sem recuo na primeira linha, justificado;
%   - frases concisas e afirmativas, verbo na voz ativa e na terceira
%     pessoa do singular; a primeira frase expõe o tema e a categoria
%     do trabalho;
%   - tipo informativo: finalidade, metodologia, resultados e conclusões;
%   - de 150 a 250 palavras no artigo;
%   - sem símbolos, fórmulas, equações ou enumeração de tópicos;
%   - espaçamento simples e tamanho 12.
\begin{center}
  \textbf{\ABNTEXchapterfont\MakeUppercase{\imprimirtitulo}}
\end{center}

% Autor e orientador em linhas distintas, alinhados à margem direita, com
% nota de rodapé de identificação (manual, seção 7.4).
\begin{flushright}
  \imprimirautor\footnote{\notadoautor}

  \imprimirorientador\footnote{\notadoorientador}

  % Havendo coorientação, remova o % da linha abaixo e crie a nota
  % correspondente em configuracoes/metadados.tex.
  % \imprimircoorientador\footnote{\notadoorientador}
\end{flushright}

\begin{center}
  \textbf{RESUMO}
\end{center}

\begingroup
\SingleSpacing
\noindent Apresente aqui, em um único parágrafo, o tema da pesquisa, seu objetivo,
os procedimentos metodológicos, os principais resultados e a conclusão.
Este texto é apenas uma orientação de preenchimento e deve ser substituído
antes da entrega. O resumo do artigo tem de 150 a 250 palavras; confira com
o curso as demais exigências.
\par
\noindent\textbf{Palavras-chave:} \palavraschave
\par
\endgroup

% ABSTRACT: escreva uma tradução fiel do resumo, não uma nova pesquisa.
% Edite as palavras-chave em inglês em configuracoes/metadados.tex.
%
% O resumo em língua estrangeira é OBRIGATÓRIO (manual IFPI 2024, seções 4
% e 7.2) e vem logo depois do resumo no idioma do documento. Pode ser em
% inglês, francês ou espanhol; para trocar o idioma, mude "english" nas
% duas linhas do ambiente otherlanguage* e o título ABSTRACT.
\begin{otherlanguage*}{english}
\begin{center}
  \textbf{ABSTRACT}
\end{center}

\begingroup
\SingleSpacing
\noindent Present the research topic, objective, methods, main results and
conclusion in a single paragraph. This is a placeholder to be replaced with
an English translation of the final Portuguese abstract before submission.
Check the word limit and other requirements with your course coordinator.
\par
\noindent\textbf{Keywords:} \keywords
\par
\endgroup
\end{otherlanguage*}

% Data de aprovação no formato dia/mês/ano, logo abaixo das keywords
% (manual IFPI 2024, seção 7.4). O valor vem de configuracoes/metadados.tex.
\begin{flushleft}
  Data de aprovação: \imprimirdataapresentacao.
\end{flushleft}


% % Ative em estrutura/pre-textual.tex. As legendas alimentam a lista.
\cleardoublepage
\pdfbookmark[0]{\listfigurename}{lof}
\listoffigures*
\cleardoublepage

% % Ative em estrutura/pre-textual.tex. As legendas alimentam a lista.
\cleardoublepage
\pdfbookmark[0]{\listtablename}{lot}
\listoftables*
\cleardoublepage

% % Ative em estrutura/pre-textual.tex. As legendas dos ambientes "quadro"
% alimentam a lista. Quadro e tabela são coisas diferentes na norma
% (manual IFPI 2024, seção 3.12): o quadro traz dados qualitativos e é
% fechado em todos os lados; a tabela traz dados numéricos e tem as
% laterais abertas. Por isso cada um tem a sua lista.
\cleardoublepage
\pdfbookmark[0]{\listadequadrosname}{loq}
\listadequadros*
\cleardoublepage

% % Ative em estrutura/pre-textual.tex. Mantenha apenas as siglas utilizadas.
\clearpage
\begin{siglas}
  \item[ABNT] Associação Brasileira de Normas Técnicas
  \item[CRB] Conselho Regional de Biblioteconomia
  \item[IBGE] Instituto Brasileiro de Geografia e Estatística
  \item[IFPI] Instituto Federal de Educação, Ciência e Tecnologia do Piauí
  \item[UFPI] Universidade Federal do Piauí
\end{siglas}

% % Ative em estrutura/pre-textual.tex. Mantenha apenas os símbolos utilizados.
%
% Manual IFPI 2024, seção 5.1.2: diferente da lista de siglas, a lista de
% símbolos segue a ORDEM EM QUE OS SÍMBOLOS APARECEM NO TEXTO, e não a
% ordem alfabética. Reordene os itens conforme o seu trabalho.
\clearpage
\begin{simbolos}
  \item [\% ] Porcentagem
  \item [ © ] Copyright
  \item [ ® ] Marca registrada
  \item [ \$ ] Dólar
  \item [ § ] Seção
  \item[$ \Gamma $] Letra grega Gama
  \item[$ \Lambda $] Lambda
  \item[$ \zeta $] Letra grega minúscula zeta
  \item[$ \in $] Pertence
\end{simbolos}


% SUMÁRIO: último elemento pré-textual, obrigatório em monografia,
% dissertação e tese (seção 5.1.2.14); não consta na estrutura do artigo.
% Para incluí-lo, ative as quatro linhas seguintes juntas.
% \cleardoublepage
% \pdfbookmark[0]{\contentsname}{toc}
% \tableofcontents*
% \cleardoublepage
